\section{Introduction}
Cardiac arrhythmias remain a leading cause of morbidity worldwide. Diagnostic capacity outside tertiary centers is constrained by access to electrophysiology expertise and the workload required for manual electrocardiogram (ECG) review. Deep neural networks now classify dozens of rhythm phenotypes from ambulatory recordings, but their adoption in telecardiology remains sensitive to performance on noisy wearable data streams and to the transparency of the generated explanations.\cite{Hannun2019Cardiologist,acharya2017deep,zheng2020residual,jenifer2024yolo} Clinicians therefore continue to request models that expose uncertainty, trace waveform rationales, and ship with complete data processing recipes.

Open cardiac signal repositories, including the MIT-BIH Arrhythmia Database, catalyze reproducible benchmarking and enable rapid comparison across architectures.\cite{Goldberger2000PhysioBank,luz2016ecg} However, many published ECG classifiers still prioritise headline accuracy over reproducibility artifacts, calibrated outputs, or diagnostics describing how minority rhythms behave. Recent biomedical engineering reports demonstrate the value of interpretable feature extraction, stringent artifact handling, and transparent documentation across biosignal modalities.\cite{rahman2024residual,rowlandson2024bridging,spagna2024telecardiology,tison2018passive} Translating these lessons back to ECG monitoring requires work that combines algorithmic detail with openly distributed preprocessing pipelines. Such pipelines allow other groups to vet device-specific adaptations and stress-test generalisation claims.

This study develops and releases a convolutional--Conformer classifier trained on Morlet wavelet representations of MIT-BIH heartbeats. It makes three contributions:
\begin{enumerate}[label=(\roman*),leftmargin=*]
  \item a lightweight neural architecture tuned through rerunnable scripts and shipped with TFRecords, split metadata, checkpoints, and plotting notebooks so reviewers and practitioners can replicate every number in the paper,
  \item comprehensive diagnostics spanning macro-averaged metrics, calibration summaries, confidence intervals, and error analyses that guide threshold selection for telecardiology workflows, and
  \item detailed documentation of preprocessing, training, and evaluation steps that aligns with emerging expectations for open methods and reproducible biosignal research.
\end{enumerate}
By focusing on transparent engineering practices and auditable baselines rather than algorithm novelty alone, the work targets the translation gap that currently limits the use of interpretable arrhythmia classifiers in routine care while offering a conservative reference point for future augmentation and cost-sensitive learning studies.
